% ------------------------------------------------------------------------------
% §6 Results — drafted Phase C from:
%   03_rq1_essentiality/findings.md + findings_3b.md; 04_rq2_compensation_price/findings.md
%   Exhibits: T1, T2, Fig 5
% ------------------------------------------------------------------------------
\section{Results}\label{sec:results}

\subsection{RQ1: essentiality moves conduct for the directed workhorse, and the
competition control absorbs none of it}

Table~\ref{tab:rq1} reports equation~\eqref{eq:rq1}. Pooled across the four
test units, the essentiality coefficient is a precise null on the fixed-\$300
definition ($-0.0045$, wild-cluster-bootstrap $p=0.85$) and negative but
insignificant on the cost-indexed definition ($-0.048$, $p=0.21$). The informative
comparison is across columns: adding the competition controls moves the
essentiality coefficient by 6.2\% and 0.7\% on the two definitions.
Whatever essentiality response exists, the energy-market competition channel
explains essentially none of it.

The pooled numbers average over opposing units. Torrens Island B --- the
station the direction regime actually leans on --- withholds \emph{more} when
essential: $-0.076$ on the cost-indexed definition ($p=0.04$), stable in sign
across all 72 leave-one-month-out re-estimates
(Section~\ref{sec:robustness}). Pelican Point's apparently positive response
($+0.18$) is a one-month artifact --- dropping October 2023 alone flips its
sign --- and is withdrawn as a finding; Table~\ref{tab:rq1} retains it,
footnoted, for completeness.\footnote{The magnitude for Torrens: $-0.076$ of
registered capacity is roughly 15~MW per 200-MW unit, on a station whose
essential-interval cheap offering is already thin.}

Two auxiliary results discipline the interpretation.

\emph{Essentiality is not scarcity.} The competition control itself shows why
the market-power reading of withholding-when-essential fails: essential
intervals are measurably \emph{more} competitive than non-essential ones
(residual-demand slope $-5.04$ versus $-3.31$ for Torrens). An
energy-market-power account of the Torrens response requires essential
periods to be tight;
they are the opposite. The alternative fails on the sign of its conditions,
not merely on magnitude.

\emph{Conduct responds to regimes, not doses.} The one strong conduct
correlate among the controls is the \emph{saturated} indicator: when rivals
are locally saturated, the unit offers 3.2 percentage points less of its
capacity cheaply ($p<0.001$), while the continuous slope carries nothing
($p=0.99$). Conduct moves with the regime of facing no local competition, not
with the degree of competition when some exists. The same shape appears on the
essentiality margin: a response to the essential state. Whether conduct also
scales with what the state \emph{pays} is exactly the RQ2 question, and under
a regime-triggered (insurance) account of withholding the answer should be no.

\subsection{RQ2: the withholding gap widens with the compensation price}

It is not no. Table~\ref{tab:rq2} reports equation~\eqref{eq:rq2}: the
essential $\times$ compensation-price interaction is $-0.051$ per \$100/MWh
on the fixed definition (analytic $p=0.005$; wild-cluster-bootstrap $p=0.004$
Rademacher, $0.004$ Webb) and $-0.055$ on the cost-indexed definition
($p=0.007$; bootstrap $0.007$ and $0.006$). The essential-versus-matched gap in the cheap
share widens by about 5.1 percentage points of registered capacity for every
\$100/MWh of compensation price --- roughly 10~MW per Torrens unit --- and by
about 13 percentage points across the observed \$121--378 range of $\dt$.

\subsection{Decomposing the margin: the response is the eligibility choice}

The share mixes two margins with different economics
(Section~\ref{sec:design}): whether the minimum-stable quantum is within
dispatch's reach at all --- the choice that averts or invites a direction
--- and how much depth is offered when it is. Table~\ref{tab:decomposition}
separates them, with interpretations committed in the project record before
estimation. The dose response lives entirely on the eligibility margin: the
probability that the floor is within reach falls by 8.6 percentage points
per \$100/MWh on essential relative to matched intervals
(wild-cluster-bootstrap $p=0.003$; randomization inference permuting the
month-to-price assignment, $p=0.046$), against a base essential-interval
reach rate of only 10\%. Conditional on the floor being offered, depth shows
no response at all ($p=0.54$--$0.79$). The share coefficient of
Table~\ref{tab:rq2} is this eligibility response diluted by unresponsive
intensive variation --- and the eligibility measure is threshold-free, since
the floor block, when offered, always sits at the $-\$1{,}000$ band. The
same gradient survives aggregation to the month grain
(Figure~\ref{fig:monthgap}): across the 19 essential-bearing months, the
monthly gap moves $-0.096$ per \$100/MWh (HC1 $p=0.001$).

\begin{table}[!htbp]\centering
\caption{Decomposing the dose response: eligibility versus intensive margin}
\label{tab:decomposition}
\begin{adjustbox}{max width=\textwidth}
\begin{threeparttable}
\begin{tabular}{lcc}
\toprule
 & Floor within reach & Offered depth, given reach \\
 & (eligibility margin) & (intensive margin) \\
\midrule
Essential $\times$ compensation price (\$100/MWh) & $-0.086^{***}$ & $0.005$ \\
 & (0.025) & (0.019) \\
Wild-cluster-bootstrap $p$ (Rademacher) & 0.003 & 0.789 \\
Range across June-2022 treatments & $-0.088$ to $-0.076$ & $0.005$ to $0.036$ \\
Randomization-inference $p$ (999 month-permutations) & 0.046 & --- \\
Month-grain WLS slope (19 months, HC1 $p$) & $-0.096$ (0.0012) & --- \\
Observations (base sample) & 140,259 & 28,555 \\
\bottomrule
\end{tabular}
\begin{tablenotes}\footnotesize
\item The eligibility margin is an indicator for the unit's minimum-stable quantum (its frozen floor: 40~MW per Torrens unit) being offered within dispatch's reach; it is identical under both cheap thresholds because the floor block, when offered, sits at the $-\$1{,}000$ band. The intensive margin is the cheap-capacity share on the subsample with the floor in reach (a decomposition aid: it conditions on an outcome). Specification, matching, controls, clustering, and June-2022 treatments exactly as in Table~\ref{tab:rq2}. The cost-indexed intensive margin behaves identically (base $
-0.019
$, $p=
0.55
$). Interpretations for all cells were committed in the project record before estimation.
\item $^{*}p<0.10$, $^{**}p<0.05$, $^{***}p<0.01$ (analytic).
\end{tablenotes}
\end{threeparttable}
\end{adjustbox}
\end{table}


\begin{figure}[!htbp]\centering
\includegraphics[width=0.85\textwidth]{fig6_month_gap_vs_price}
\caption{The month-grain version of the dose response: the monthly
essential-vs-matched gap in floor-reach against the compensation price,
19 months with at least 30 essential unit-rows, point size proportional to
essential mass. The fitted line is the essential-mass-weighted least-squares
slope ($-0.096$ per \$100/MWh, HC1 $p=0.001$).}
\label{fig:monthgap}
\end{figure}

Unlike the RQ1 level result, the interaction owes nothing to June 2022:
across all four treatments of the crisis month the estimate sits between
$-0.044$ and $-0.058$ with $p$-values of 0.001 to 0.029
(Figure~\ref{fig:rq2stability}). The essential \emph{main effect} in
Table~\ref{tab:rq2} is positive ($+0.089$): at low compensation prices,
essential intervals show \emph{more} cheap capacity than matched ones ---
these are winter-tight periods in which running pays --- and the gap
crosses zero at $\dt\approx\$175$, visible directly in the month-grain
scatter (Figure~\ref{fig:monthgap}). The dose response is not a deepening
of an always-negative gap; it is a reversal, from relative presence when
the direction channel pays little to relative absence as it pays more. Attenuation from realised-state
classification biases $\theta$ toward zero, so the true payment-sensitivity is
if anything larger.

The channel ranking, under the interpretation committed before estimation,
is then: energy-market power --- eliminated on the wrong sign of its
conditions; pure presence-inelasticity with no payment sensitivity --- the
insurance account of Section~\ref{sec:design} --- rejected, because it
predicted a null here; payment-seeking --- the account most consistent with
the data. One refinement of the regime-not-dose pattern survives
alongside: conduct is regime-triggered on the \emph{competition} margin
(saturation, not slope), while on the \emph{payment} margin the essential-state
response scales with the prize.

Two limitations are reported with the estimate, as committed with the
design before estimation. The interaction is identified from cross-month
variation; anything that widens the essential-versus-matched gap \emph{and}
co-moves with the compensation price across months --- 2022H2 fuel-supply
stress is the obvious candidate --- is not excluded by month effects.
Section~\ref{sec:robustness} confronts this directly with the timing wedge
the formula creates (gas fell nine months before the payment did) and a
fuel horse race; the pure fuel story fails its point predictions there, but
the sharpest formula-driven contrast is marginal, so this paper's causal
language for the cross-month gradient is deliberately ``consistent with
payment-seeking,'' not proof. And the essential mass is concentrated: 21
effective month-clusters, half the mass in three of them. The estimate
clears conventional significance under bootstrap inference designed for
this cluster structure and under randomization inference
(Table~\ref{tab:decomposition}), but the concentration stands as a stated
limitation.

A third limitation deserves its own statement: the treatment is not
mechanically exogenous to the units that receive it. The compensation price
is a percentile of past regional spot prices, and withholding that raises
spot prices feeds, with a lag, the units' own future payment. Three facts
bound the loop without eliminating it. Contemporaneously it is negligible
by construction --- a month's own conduct enters its own $\dt$ only through
that month's marginal contribution to a 365-day window. Dynamically, the
loop requires the units' absence to move the \emph{top-decile} prices that
set the percentile, and the record places their absence where prices are
low: essential intervals are the market's more-competitive, cheap-energy
periods (Section~\ref{sec:results}), not its price spikes. And the one
large movement in $\dt$ over the sample --- the 2022 surge to \$349 ---
tracked NEM-wide fuel costs (\$322/MWh at the focal units' own heat rates),
a system event no three units authored. What survives as a caveat is the
softer form: any residual own-price influence inflates the payment
gradient's causal reading, and the paper's ``consistent with'' language is
calibrated to that residual as well.

\begin{table}[!htbp]\centering
\caption{RQ1: Essentiality and the cheap-capacity share}
\label{tab:rq1}
\begin{adjustbox}{max width=\textwidth}
\begin{threeparttable}
\begin{tabular}{lcccccc}
\toprule
 & \multicolumn{2}{c}{M1} & \multicolumn{2}{c}{M2} & \multicolumn{2}{c}{M3} \\
\cmidrule(lr){2-3}\cmidrule(lr){4-5}\cmidrule(lr){6-7}
 & Fixed \$300 & 2$\times$SRMC & Fixed \$300 & 2$\times$SRMC & Fixed \$300 & 2$\times$SRMC \\
\midrule
\multicolumn{7}{l}{\emph{Panel A: pooled (four focal units)}} \\
Essential & $-0.004$ & $-0.048$ & $-0.004$ & $-0.048$ & $-0.004$ & $-0.048$ \\
 & (0.023) & (0.038) & (0.023) & (0.038) & (0.023) & (0.038) \\
 & [0.847] & [0.208] & [0.847] & [0.207] & [0.857] & [0.207] \\
Saturated (slope $=0$) & & & & & $-0.032^{***}$ & $-0.036^{***}$ \\
 & & & & & (0.006) & (0.007) \\
\midrule
\multicolumn{7}{l}{\emph{Panel B: heterogeneity (M3 specification)}} \\
Essential, Torrens Island B only & & & & & $-0.037$ & $-0.076^{**}$ \\
 & & & & & (0.023) & (0.036) \\
Essential, Pelican Point only\tnote{a} & & & & & $0.184^{*}$ & $0.128$ \\
 & & & & & (0.099) & (0.084) \\
\midrule
Observations (pooled) & \multicolumn{6}{c}{1,261,576} \\
\bottomrule
\end{tabular}
\begin{tablenotes}\footnotesize
\item Outcome: share of registered capacity offered below the cheap threshold (fixed \$300/MWh or 2$\times$ engineering SRMC), 5-minute unit--interval panel, 2022--2024. M1 adds demand, non-synchronous supply, SRMC, and price controls; M2 adds the residual-demand slope; M3 adds the saturated (zero-slope) indicator. Analytic standard errors clustered by month in parentheses; wild cluster bootstrap $p$-values (Rademacher, $R=999$, 35 df) in brackets for the essentiality coefficient. Panel B re-estimates M3 on the Torrens Island B and Pelican Point subsamples.
\item[a] The Pelican Point coefficient is not robust: it flips sign when October 2023 is dropped (leave-one-month-out, Section~\ref{sec:robustness}) and is withdrawn as a finding.
\item $^{*}p<0.10$, $^{**}p<0.05$, $^{***}p<0.01$ (analytic).
\end{tablenotes}
\end{threeparttable}
\end{adjustbox}
\end{table}

\begin{table}[!htbp]\centering
\caption{RQ2: The essential-interval withholding gap widens with the compensation price}
\label{tab:rq2}
\begin{adjustbox}{max width=\textwidth}
\begin{threeparttable}
\begin{tabular}{lccc}
\toprule
 & \multicolumn{2}{c}{Essential $\times$ compensation price (\$100/MWh)} & \\
\cmidrule(lr){2-3}
Sample treatment of June 2022 & Fixed \$300 & 2$\times$SRMC & $N$ \\
\midrule
Base (exclude June-2022 suspension window) & $-0.051^{***}$ & $-0.055^{***}$ & 140,259 \\
 & (0.016) & (0.019) & \\
(i) Exclude all June 2022 & $-0.051^{***}$ & $-0.058^{***}$ & 134,636 \\
 & (0.016) & (0.016) & \\
(ii) Include window at APC \$300 & $-0.044^{**}$ & $-0.050^{**}$ & 144,358 \\
 & (0.019) & (0.020) & \\
(iii) Base minus pre-suspension June & $-0.051^{***}$ & $-0.058^{***}$ & 135,269 \\
 & (0.016) & (0.016) & \\
\midrule
Essential (main effect, base sample) & $0.089^{*}$ & $0.083$ & \\
 & (0.049) & (0.058) & \\
\midrule
WCB $p$ (base row): Rademacher & 0.004 & 0.007 & \\
\phantom{WCB $p$ (base row):} Webb & 0.004 & 0.006 & \\
\bottomrule
\end{tabular}
\begin{tablenotes}\footnotesize
\item Coefficient on essential $\times$ compensation price (per \$100/MWh) on the CEM-matched sample (unit $\times$ month $\times$ non-synchronous quintile $\times$ hour block $\times$ competition bin). A coefficient of $-0.051$ means the essential-vs-matched gap in the cheap-capacity share widens by 5.1 percentage points of registered capacity per \$100/MWh. Analytic standard errors clustered by month in parentheses; specification and controls as in Table~\ref{tab:rq1} (M3). Wild cluster bootstrap $p$-values ($R=999$, 20 df) for the base sample. The interpretation of this test, including the reading of a null, was committed in the project record before estimation.
\item $^{*}p<0.10$, $^{**}p<0.05$, $^{***}p<0.01$ (analytic).
\end{tablenotes}
\end{threeparttable}
\end{adjustbox}
\end{table}


\begin{figure}[!htbp]\centering
\includegraphics[width=0.85\textwidth]{fig5_rq2_stability}
\caption{The RQ2 interaction across June-2022 treatments and outcome
definitions. Points are estimates of $\theta$ in
equation~\eqref{eq:rq2}; bars are 95\% analytic confidence intervals,
clustered by month. All eight estimates lie between $-0.044$ and $-0.058$ per
\$100/MWh.}
\label{fig:rq2stability}
\end{figure}
